\documentclass[10pt,a4paper]{amsart}
\usepackage[margin=19mm]{geometry}
\usepackage{amsmath,amssymb,mathtools}
\usepackage[scr=rsfs,cal=euler]{mathalfa}
\usepackage{xfrac}
\usepackage[unicode,hidelinks,bookmarks=false]{hyperref}
\newcommand{\ie}{i.e.\ }
\newcommand{\cf}{cf.\ }
\newcommand{\resp}{respectively\ }
\newcommand{\Subsection}{\subsection}
\providecommand{\nobreakdash}{\nobreak\hbox{-}}
\setlength{\emergencystretch}{2em}
\multlinegap0pt
\usepackage{bm}
\usepackage{bbm}
\usepackage{dsfont}
\usepackage{xspace}
\usepackage{cancel}
\usepackage{mathtools}
\usepackage{amsthm}
\makeatletter
\@ifundefined{lemma}{\newtheorem{lemma}{Lemma}}{}
\makeatother
\newcommand{\A}{\mathcal{A}}
\newcommand{\N}{\mathbb{N}}
\newcommand{\var}{\operatorname{var}}
\title[Central Limit Theorem for tensor products of free variables]{Central Limit Theorem for tensor products of free variables}
\author{C\'ecilia Lancien, Patrick Oliveira Santos, Pierre Youssef}
\date{Source: arXiv:2404.19662v1 (CC BY 4.0)}
\begin{document}
\maketitle

\noindent\emph{Source and licence.} Central Limit Theorem for tensor products of free variables. Version arXiv:2404.19662v1 (CC BY 4.0), distributed under the Creative Commons Attribution 4.0 International licence, as recorded in the arXiv licence metadata for this exact version and shown on the abstract page https://arxiv.org/abs/2404.19662v1. Full rights notice: licenses/arxiv-2404.19662-CC-BY-4.0.txt.\par\smallskip

\noindent\textit{Editorial scope.} Counted unit: the Lemma whose source label is \texttt{lemma: case theta\_1=1} in the article (source0.tex lines 1007--1021, source label \texttt{lemma: case theta\_1=1}), delivered together with the article's own complete original proof of it (source0.tex lines 1022--1147, one \texttt{proof} environment of 126 lines). No mathematics is written, repaired or re-derived: statement and proof are the article's own text line for line. Required context retained uncounted: lemma: contribution of ac1ac2 (lines 448--453); lemma: centering (lines 643--656); equation: a\_j is either a or tilde a (lines 907--909); eq: power of delta (lines 1001--1003). Every definition, display and label of those lines is retained unchanged. Omitted from this package and not redistributed: 1--447, 454--642, 657--906, 910--1000, 1004--1006, 1148--1534. Nothing inside a retained excerpt is omitted or altered: every retained body is delivered exactly as the article's lines are written, so no omission receipt of any kind accompanies this package. Citations: the retained excerpts carry no citation command at all, so no bibliography entry of the article is transcribed and no reference is redistributed. Reference adaptation: none was needed and none was made; every reference inside the delivered unit resolves inside it, so no unresolved reference or citation is produced. Printed slips kept literal: no printed word, symbol, hypothesis or number of the counted statement or of its proof was altered. Layout and class: the article's own class is \texttt{amsart} with packages including fontenc, babel, inputenc, natbib, amssymb, amsmath, xcolor, amsfonts, amsthm, csquotes, geometry, graphicx, comment, tikz; the wrapper is the lane's pinned \texttt{amsart} editorial wrapper at 10pt on A4, which restates the theorem-like environments the retained text needs and adds the article's own macro definitions that the retained text needs (\texttt{\textbackslash A}, \texttt{\textbackslash N}, \texttt{\textbackslash var}), with extra wrapper packages (bm, bbm, dsfont, xspace, cancel, mathtools). The delivered numbering is the unit's own. Assets: no retained span contains a figure environment, a graphics-inclusion command, a PDF-page inclusion, a table or a TikZ picture; no image file is copied into this package, the package declares no asset dependency and no image file is redistributed with it. Licence: version arXiv:2404.19662v1 of this article is distributed under the Creative Commons Attribution 4.0 International licence, as recorded in the arXiv licence metadata for this exact version and shown on the abstract page https://arxiv.org/abs/2404.19662v1. A full rights notice is delivered at licenses/arxiv-2404.19662-CC-BY-4.0.txt. Characters: every retained excerpt is pure ASCII, so the delivered engine, XeLaTeX with its default Latin Modern text font, reports no missing character.
\begin{lemma}\label{lemma: contribution of ac1ac2}
    Let $a,c_1,c_2$ be variables such that $a$ is free from $\{c_1,c_2\}$. Then
    \begin{align*}
        \tau(ac_1ac_2)=\var(a)\tau(c_1)\tau(c_2)+\tau^2(a)\tau(c_1c_2).
    \end{align*}
\end{lemma}

\begin{lemma}\label{lemma: centering}
    Let $(\A,\tau)$ be a unital faithful tracial noncommutative probability space. Let $\{a_k: k \in [n]\},\{d_k: k \in [n]\} \subset \A$ be two collections of free random variables and let
    \begin{align*}
        c_k:=a_k\otimes d_k-\tau\otimes \tau(a_k\otimes d_k),
    \end{align*}
    for every $k\in[n]$. Then for any $m \ge 1$ and $i_1,\ldots,i_m \in [n]$ such that there exists an index $l^* \in [m]$ satisfying $i_{l^*} \ne i_j$ for all $j \ne l^*$, we have
    \begin{align*}
        \tau\otimes \tau(c_{i_1}\cdots c_{i_{m}})=0.
    \end{align*}
    In particular, we have
    \begin{align*}
        \tau\otimes \tau\left(c_{i_1}\cdots c_{i_{l^*-1}}a_{i_{l^*}}\otimes d_{i_{l^*}} c_{i_{l^*+1}}\cdots c_{i_m} \right)=\tau(a_{i_{l^*}})\tau(d_{i_{l^*}})\tau\otimes \tau(c_{i_1}\cdots c_{i_{l^*-1}} c_{i_{l^*+1}}\cdots c_{i_m}).
    \end{align*}
\end{lemma}

\begin{align}\label{equation: a_j is either a or tilde a}
    a_{i_j}^{(l)},a_{i_j}^{(r)} \in \{a_{i_j},\tilde{a}_{i_j}\},
\end{align}

\begin{align}\label{eq: power of delta}
    h:=|I_1|+|I_2|+|I_3|=p-2(k+1).
\end{align}

\begin{lemma}[Case $\theta_1=1$]\label{lemma: case theta_1=1}
    Let $p \ge 4$ be an even integer, $\pi \in P_2^{\operatorname{con}}(p)$. Let $V_1,\ldots,V_k \in \pi$ be blocks, $w \in \{0,1\}^k$ be a coloring of $V_1,\ldots,V_k$ and
    \begin{align*}
        B_k=B_k\left(V_1^{(w_1)},\ldots,V_k^{(w_k)}\right)=\left(b_{i_j,k}\right)_{j \in [p]}.
    \end{align*}
    Let $V=\{v_1<v_2\} \in \pi\setminus \{V_1,\ldots,V_{k}\}$ and the joint law 
    \begin{align*}
        \left(b_{i_{v_1},k},b_{i_{v_2},k}\right)\overset{d}{=}\left(b_{i_{v_1}}(\theta),b_{i_{v_2}}(\theta)\right).
    \end{align*}
    If $\theta_1=1$, we have
    \begin{align*}
        &\tau\otimes\tau \left(\pi\setminus \{V_1,\ldots,V_{k}\},V_1^{(w_1)},\ldots,V_{k}^{(w_{k})}\right)\\
        &=\frac{q}{2}\sum_{w=0,1}\tau\otimes\tau \left(\pi\setminus \{V_1,\ldots,V_k,V\},V_1^{(w_1)},\ldots,V_{k}^{(w_{k})},V^{(w)}\right).
    \end{align*}
\end{lemma}

\begin{proof}
    To simplify the notation, let
    \begin{align*}
        T:=\tau\otimes\tau \left(\pi\setminus \{V_1,\ldots,V_{k}\},V_1^{(w_1)},\ldots,V_{k}^{(w_{k})}\right).
    \end{align*}
    Then, we must prove that
    \begin{align*}
        T=\frac{q}{2}\sum_{w=0,1}\tau\otimes\tau \left(\pi\setminus \{V_1,\ldots,V_k,V\},V_1^{(w_1)},\ldots,V_{k}^{(w_{k})},V^{(w)}\right).
    \end{align*}
    As $\theta_1=1$, we can write
    \begin{align*}
        &b_{i_{v_1},k}=\frac{1}{\delta}\left(a\otimes a-\lambda^2\right);\\
        &b_{i_{v_2},k}=\frac{1}{\delta}\left(a\otimes a-\lambda^2\right),
    \end{align*}
    where $a:=a_{i_{v_1}}$. We get by Lemma \ref{lemma: centering} that
    \begin{align}\label{eq: value of T}
        T&=\frac{1}{\delta^2}\tau\otimes \tau\left\{\left(B_k\right)_{I_1}a\otimes a\left(B_k\right)_{I_2}a\otimes a\left(B_k\right)_{I_3}\right\}\nonumber\\
    & \quad -\frac{\lambda^4}{\delta^2}\tau\otimes \tau\left(\pi\setminus\{V_1,\ldots,V_k,V\},V_1^{(w_1)},\ldots,V_k^{(w_k)}\right).
    \end{align}
    Let us write 
    \begin{align*}
        B_k=\frac{1}{\delta}\left(a_{i_j}^{(l)}\otimes a_{i_j}^{(r)}-\lambda^2\right)_{j \in [p]}.
    \end{align*}
    Then
    \begin{align*}
        &\tau\otimes \tau\left\{\left(B_k\right)_{I_1}a\otimes a\left(B_k\right)_{I_2}a\otimes a\left(B_k\right)_{I_3}\right\}\\
        &=\frac{1}{\delta^{h}}\sum_{\substack{J_1 \subseteq I_1\\ J_2 \subseteq I_2\\ J_3 \subseteq I_3}}(-1)^{|J_1|+|J_2|+|J_3|}\lambda^{2(|J_1|+|J_2|+|J_3|)}\tau\left(a_{J_1^c}^{(l)}\, a \, a_{J_2^c}^{(l)}\, a \, a_{J_3^c}^{(l)}\right)\tau\left(a_{J_1^c}^{(r)}\, a \, a_{J_2^c}^{(r)}\, a\, a_{J_3^c}^{(r)}\right),
    \end{align*}
    where $h$ is defined as in \eqref{eq: power of delta}. We use Lemma \ref{lemma: contribution of ac1ac2} and compute
    \begin{align*}
        &R(J_1,J_2,J_3) \\
        &:=\tau\left(a_{J_1^c}^{(l)}\, a \, a_{J_2^c}^{(l)}\, a \, a_{J_3^c}^{(l)}\right)\tau\left(a_{J_1^c}^{(r)}\, a \, a_{J_2^c}^{(r)}\, a\, a_{J_3^c}^{(r)}\right)\\
        &=\left(\sigma^2\tau(a_{J_1^c}^{(l)}\, a_{J_3^c}^{(l)})\tau(a_{J_2^c}^{(l)})+\lambda^2\tau(a_{J_1^c}^{(l)}\, a_{J_2^c}^{(l)}\, a_{J_3^c}^{(l)})\right)
        \left(\sigma^2\tau(a_{J_1^c}^{(r)}\, a_{J_3^c}^{(r)})\tau(a_{J_2^c}^{(r)})+\lambda^2\tau(a_{J_1^c}^{(r)}\, a_{J_2^c}^{(r)}\, a_{J_3^c}^{(r)})\right).
    \end{align*}
    After the expansion, we have
    \begin{align*}
        R(J_1,J_2,J_3)=:\sigma^4R_1(J_1,J_2,J_3)+\sigma^2\lambda^2 \left(R_2(J_1,J_2,J_3)+R_3(J_1,J_2,J_3)\right)+\lambda^4R_4(J_1,J_2,J_3),
    \end{align*}
    where we choose $R_2$ to be the term with factor $\tau(a_{J_1^c}^{(l)}\, a_{J_2^c}^{(l)}\, a_{J_3^c}^{(l)})$. For $R_1$, we have
    \begin{align*}
        &\frac{1}{\delta^h}\sum_{\substack{J_1 \subseteq I_1\\ J_2 \subseteq I_2\\ J_3 \subseteq I_3}}(-1)^{|J_1|+|J_2|+|J_3|}\lambda^{2(|J_1|+|J_2|+|J_3|)}R_1(J_1,J_2,J_3)\\
        &=\tau\otimes \tau \left(\left(B_k\right)_{I_1}\left(B_k\right)_{I_3}\right)\tau\otimes \tau\left(\left(B_k\right)_{I_2}\right).
    \end{align*}
    Since $\pi$ is connected, $V$ crosses at least one block $U=\{j,j'\} \in \pi$. Therefore, we assume $j \in I_2$ whose matching symbol $j' \in I_1\cup I_3$. By Lemma \ref{lemma: centering}, we have
    \begin{align*}
        \tau\otimes \tau\left(\left(B_k\right)_{I_2}\right)=0,
    \end{align*}
    and the contribution of $R_1$ is zero. For $R_4$, we have
    \begin{align*}
        &\frac{1}{\delta^h}\sum_{\substack{J_1 \subseteq I_1\\ J_2 \subseteq I_2\\ J_3 \subseteq I_3}}(-1)^{|J_1|+|J_2|+|J_3|}\lambda^{2(|J_1|+|J_2|+|J_3|)}R_4(J_1,J_2,J_3)\\
        &=\tau\otimes \tau \left(\left(B_k\right)_{I_1}\left(B_k\right)_{I_2}\left(B_k\right)_{I_3}\right).
    \end{align*}
    Since
    \begin{align*}
        \tau\otimes \tau \left(\left(B_k\right)_{I_1}\left(B_k\right)_{I_2}\left(B_k\right)_{I_3}\right)=\tau\otimes\tau\left(\pi\setminus\{V_1,\ldots,V_k,V\},V_1^{(w_1)},\ldots,V_k^{(w_k)}\right),
    \end{align*}
    its contribution cancels out with $\frac{\lambda^4}{\delta^2}\tau\otimes \tau\left(\pi\setminus\{V_1,\ldots,V_k,V\},V_1^{(w_1)},\ldots,V_k^{(w_k)}\right)$ in \eqref{eq: value of T}. For $R_2$, let us first denote
    \begin{align*}
        B_{k+1}=B_{k+1}\left(V_1^{(w_1)},\ldots,V_{k}^{(w_k)},V^{(1)}\right)=:\frac{1}{\delta}\left(\overline{a}_{i_j}^{(l)}\otimes \overline{a}_{i_j}^{(r)}-\lambda^2\right)_{j \in [p]}.
    \end{align*}
    We then have
    \begin{align*}
        &\tau\otimes\tau\left(\pi\setminus\{V_1,\ldots,V_k,V\},V_1^{(w_1)},\ldots,V_{k}^{(w_k)},V^{(1)}\right)\\
        &=\frac{1}{\delta^h}\sum_{\substack{J_1 \subseteq I_1\\ J_2 \subseteq I_2\\ J_3 \subseteq I_3}}(-1)^{|J_1|+|J_2|+|J_3|}\lambda^{2(|J_1|+|J_2|+|J_3|)}\tau\left(\overline{a}_{J_1^c}^{(l)}\, \overline{a}_{J_2^c}^{(l)}\, \overline{a}_{J_3^c}^{(l)}\right)\tau\left(\overline{a}_{J_1^c}^{(r)}\, \overline{a}_{J_2^c}^{(r)}\, \overline{a}_{J_3^c}^{(r)}\right).
    \end{align*}
    By definition, the color $w_{k+1}=1$ does not affect the left legs, hence
    \begin{align*}
        \overline{a}_{i_j}^{(l)}=a_{i_j}^{(l)}.
    \end{align*}
    For the right legs and a block $U=\{t_1<t_2\}$, let us divide it into two cases.
    \begin{enumerate}
        \item\label{case 1} If $U$ does not cross $V$, we have
        \begin{align*}
            &\overline{a}_{i_{t_1}}^{(r)}=a_{i_{t_1}}^{(r)};\\
            &\overline{a}_{i_{t_2}}^{(r)}=a_{i_{t_2}}^{(r)}.
        \end{align*}
        \item\label{case 2} Otherwise, either $t_1 \in I_2$ and $t_2 \in I_1\cup I_3$ or $t_1 \in I_1 \cup I_3$ and $t_2 \in I_2$. By the definition of the color $w_{k+1}=1$, we have
        \begin{align*}
            &\overline{a}_{i_{t_1}}^{(r)}=\tilde{a}_{i_{t_1}};\\
            &\overline{a}_{i_{t_2}}^{(r)}=a_{i_{t_1}}.
        \end{align*}
    \end{enumerate}
    Since $(a_k)_{k \in \N}$ and $(\tilde{a}_k)_{k \in \N}$ are free i.i.d copies of $a$, the following holds. First, let $t \in I_2$. Consider $t'$ its matching symbol, that is, $\{t,t'\} \in \pi$ and $i_t=i_{t'} \ne i_j$, for all $j \notin \{t,t'\}$. Recall that, by the definition in \eqref{equation: a_j is either a or tilde a}, we have
    \begin{align*}
        \overline{a}_{i_t}^{(r)},\overline{a}_{i_{t'}}^{(r)} \in \{a_{i_t},\tilde{a}_{i_t}\}.
    \end{align*}
    If $t' \in I_2$, Case \eqref{case 1} implies that the law of $(\overline{a}_{i_t}^{(r)},\overline{a}_{i_{t'}}^{(r)})$ is equal to the law of $({a}_{i_t}^{(r)},{a}_{i_{t'}}^{(r)})$. If $t' \in I_1\cup I_3$, as $a_{i_t},\tilde{a}_{i_t}$ are free i.i.d and free from $(a_k)_{k \ne i_t},(\tilde{a}_k)_{k \ne i_t}$, we have
    \begin{align}\label{eq: equality in distribution t' not in I_2}
        \left((\overline{a}_{i_j}^{(r)})_{j \in I_2\setminus\{t\}},\overline{a}_{i_t}^{(r)}\right)\overset{d}{=}\left((\overline{a}^{(r)}_{i_j})_{j \in I_2\setminus\{t\}},a_{i_t}\right)\overset{d}{=}\left((\overline{a}^{(r)}_{i_j})_{j \in I_2\setminus\{t\}},a_{i_t}^{(r)}\right).
    \end{align}
    Applying the above equalities in distribution for all $t \in I_2$, namely, Case \eqref{case 1} for $t' \in I_2$ and \eqref{eq: equality in distribution t' not in I_2} for $t' \notin I_2$, we get
    \begin{align*}
        \left(\overline{a}_{i_j}^{(r)}\right)_{j \in I_2}\overset{d}{=}\left(a_{i_j}^{(r)}\right)_{j \in I_2}.
    \end{align*}
    By symmetry, we also have
    \begin{align*}
        \left(\overline{a}_{i_j}^{(r)}\right)_{j \in I_1\cup I_3}\overset{d}{=}\left(a_{i_j}^{(r)}\right)_{j \in I_1\cup I_3}.
    \end{align*}
    Additionally, $\left(\overline{a}_{i_j}^{(r)}\right)_{j \in I_2}$ is now free from $\left(\overline{a}_{i_j}^{(r)}\right)_{j \in I_1\cup I_3}$. Indeed, let us show that each variable $\overline{a}_{i_t}^{(r)}$ is free from $\left(\overline{a}_{i_j}^{(r)}\right)_{j \in I_1\cup I_3}$, for all $t \in I_2$. Consider $t'$ its matching symbol. First, if $t' \in I_2$, the fact that $(a_k)_{k \in \N}$ and $(\tilde{a}_k)_{k \in \N}$ are free i.i.d collections implies that $\overline{a}_{i_t}^{(r)}$ (and $\overline{a}_{i_{t'}}^{(r)}$) is free from $\left(\overline{a}_{i_j}^{(r)}\right)_{j \in I_1\cup I_3}$. If $t' \in I_1\cup I_3$, Case \eqref{case 2} implies that $\overline{a}_{i_t}^{(r)}$ is free from $\overline{a}_{i_{t'}}^{(r)}$ and then again $\overline{a}_{i_t}^{(r)}$ is free from $\left(\overline{a}_{i_j}^{(r)}\right)_{j \in I_1\cup I_3}$. 
    
    In summary, by freeness, we have
    \begin{align*}
        \tau\left(\overline{a}_{J_1^c}^{(l)}\, \overline{a}_{J_2^c}^{(l)}\, \overline{a}_{J_3^c}^{(l)}\right)\tau\left(\overline{a}_{J_1^c}^{(r)}\, \overline{a}_{J_2^c}^{(r)}\, \overline{a}_{J_3^c}^{(r)}\right)&=\tau\left(a_{J_1^c}^{(l)}\, a_{J_2^c}^{(l)}\, a_{J_3^c}^{(l)}\right)\tau\left(a_{J_1^c}^{(r)}\, a_{J_3^c}^{(r)}\right)\tau\left(a_{J_2^c}^{(r)}\right)\\
        &=R_2(J_1,J_2,J_3).
    \end{align*}
    Therefore
    \begin{align*}
        &\frac{1}{\delta^h}\sum_{\substack{J_1 \subseteq I_1\\ J_2 \subseteq I_2\\ J_3 \subseteq I_3}}(-1)^{|J_1|+|J_2|+|J_3|}\lambda^{2(|J_1|+|J_2|+|J_3|)}R_2(J_1,J_2,J_3)\\
        &=\tau\otimes \tau\left(\pi\setminus\{V_1,\ldots,V_k,V\},V_1^{(w_1)},\ldots,V_k^{(w_k)},V^{(1)}\right),
    \end{align*}
    Following the same reasoning, it follows that
    \begin{align*}
        &\frac{1}{\delta^h}\sum_{\substack{J_1 \subseteq I_1\\ J_2 \subseteq I_2\\ J_3 \subseteq I_3}}(-1)^{|J_1|+|J_2|+|J_3|}\lambda^{2(|J_1|+|J_2|+|J_3|)}R_3(J_1,J_2,J_3)\\
        &=\tau\otimes \tau\left(\pi\setminus\{V_1,\ldots,V_k,V\},V_1^{(w_1)},\ldots,V_k^{(w_k)},V^{(0)}\right).
    \end{align*}
    Combining the above relations in $T$, we get
    \begin{align*}
        T=\frac{\sigma^2\lambda^2}{\delta^2}\sum_{w=0,1}\tau\otimes \tau\left(\pi\setminus\{V_1,\ldots,V_k,V\},V_1^{(w_1)},\ldots,V_k^{(w_k)},V^{(w)}\right).
    \end{align*}
    It remains to note that 
    \begin{align*}
        \frac{\sigma^2\lambda^2}{\delta^2}=\frac{q}{2},
    \end{align*}
    to finish the proof. 
\end{proof}

\end{document}
